\documentclass[11pt]{article}
\usepackage{amsmath}
\usepackage{graphicx}
\usepackage{float}
\usepackage{blindtext}
\usepackage{subcaption}
\usepackage[a4paper,margin=2cm]{geometry}
\usepackage{fancyhdr}
\usepackage{longtable}
\usepackage[
    backend=biber,
    style=ieee
]{biblatex}
\usepackage{xcolor}
\usepackage{soul}
\usepackage{listings}

\newcommand{\hlc}[2][yellow]{{%
    \colorlet{foo}{#1}%
    \sethlcolor{foo}\hl{#2}}%
}


\addbibresource{report.bib}

\graphicspath{ {./images/} }

\pagestyle{fancy}
\fancyhead[C]{Report}
\fancyhead[L]{EEE3027 Calculator Assignment}
\fancyhead[R]{2nd of May}

\cfoot{} % get rid of the page number 
\fancyfoot[L]{6596386}
\fancyhead[C]{}
\fancyfoot[R]{\thepage}

\begin{document}
\begin{titlepage}
    \begin{center}
    \includegraphics[width=\textwidth]{Logo.png} % also works with logo.pdf
    \vfill
    \Huge
    \textbf{EEE3027: Calculator Assignment}
    \vfill
    \huge
    Assignment Report\\
    \vspace{1cm}
    \Large
    2nd of May 2023\\
    URN: 6596386\\
    \vfill
    \vfill
    \Large
    Department of Electronic Engineering\\
    Faculty of Engineering and Physical Sciences\\
    University of Surrey\\
    \end{center}
\end{titlepage}

\begin{abstract}
This report covers the work done on creating and testing an FPGA Calculator in VHDL as part of EEE3027 Digital Design with VHDL.
VHDL was used to program an FPGA capable of taking serial input,
converting the serial input to numbers and operations,
calculating the result,
converting the result back into serial, 
and finally transmitting the result back.
The final design allows for a user to input, over serial, any integer between -32767 and 32767 (a signed 16-bit value) with the following integer operation : addition, subtraction, multiplication, and division. 
The FPGA will then return the result if it can be represented as a signed 16-bit value, otherwise it will return an error. 
The report also covers background theory for UART communication, an exploration of the original provided code, and a discussion on timing and physical implementation on a Minized Zynq board.
\vspace{1cm}
\end{abstract}
\tableofcontents
\pagebreak

\section{Overview}

The objectives of this assignment were to create an arithmetic calculator using the std.numeric library and provided UART code.
This provided code first has to be debugged and tested.
The calculators basic requirement was to compute values given in the "$A {+-*/} B$" format, with no specific instructions of the range of values of $A$, $B$ or the output.
For the implementation created for this assignment the calculator is able to take inputs and provide outputs that can be described as signed 16-bit integers (ie between -32767 and 32767),
however, through the use of generics this range can be changed.
The serial UART input is then processed into the two expected numbers and operations, with any character that is not a numeral (0 to 9) or an operator character ($+-*/$) being ignored.
At the end of the input a carriage return character (ASCII Hex code $0D$) is required to indicate the end of number $B$.
A carriage return is also transmitted at the end of the calculator response with the result.
Furthermore the final design will output a special character, "!", when an operation is invalid (such as dividing by zero) or the output is outside the calculators range.

\begin{figure}[H]        
    \centering
    \includegraphics[width=\textwidth]{Logo.png}
    \caption{Top Level Implementation of Calculator}
    \label{fig:toplevel}
\end{figure} 

This achieve this functionality the design seen in figure \ref{fig:toplevel} was created, with four main components all interacting together.
These components are the UART block, the ASCII decoder, the calculator, ASCII encoder.
The UART block handles the receiving of serial UART data and converting it to parallel data as well as taking parallel data and sending it as serial UART.
The ASCII decoder takes the parallel data, assuming they are ASCII codes, and converts into two numbers and a operation.
The calculator computes operations between two numbers, the provides a result as well as a signal if the operation is invalid.
THe ASCII encoder takes a result, as well the invalid signal, and converts them into a sequences of parallel ASCII codes that it sends to the UART block.
Details on the design and operation of all of these components, as well information on testbenches and implementation on a Minized Zynq FPGA board are provided in


\section{UART}
\subsection{UART Theory}
UART stands for universal asynchronous receiver-transmitter and is a very common device-to-to device communication protocol\cite{UART}.
For this assignment a two-way UART connection is utilized,
allowing two devices (the FPGA and a connected PC) to transfer data between each other regardless of clock rate as long as they communicate over a predefined rate know as the baud rate.
Unlike other communication protocols such as SPI or I2C, UART only ever has one "master" device and one "slave" device a third device cannot be added on the same lines.

UART takes parallel data and converts it into a bit by bit frame to send.
The frame always start with the start bit, which takes the line low for one cycle as UART is normally high.
then each data bit is sent one by one (least significant bit first) with a 1 being high and a zero being low.
Some implementation will include a parity bit after the data, which will be 0 if the data total is odd and 1 is the total is even.
Parity is useful to ensure that data been altered during transmission, which can occur due to "electromagnetic radiation, mismatched baud rates, or long-distance data transfers"\cite{UART}.
The provided UART however does not utilize the parity bit.
At the end of the message a stop bit is sent going from low to high.
Equation \ref{eq:uart_sample} shows an example of how data gets converted to 

\begin{align*}
    \text{Data} &= 01000001
    \stepcounter{equation}\tag{\theequation}\label{eq:uart_sample} \\
    \text{UART Frame} &=    \overset{\text{Start Bit}}{\{0\}}
                            \underset{\text{Data}}{\{1000 0010\}}
                            \overset{\text{Stop Bit}}{\{1\}}\\
\end{align*}


The final key concept behind UART is the baud rate, which the rate of which information is transferred.
In the case of UART baud rate represents how long each bit is held high or low.
As UART is asynchronous, each device needs to know baud rate ahead of time and generate its own clock at the baud rate.
This is enough to transmit UART as the line can simply be updated each clock cycle,
however for receiving the line needs to be checked more often as the real line will often have non-instantaneous transition periods from high to low meaning that is best to check the signal in the middle of each cycle.
To do this there is often a secondary oversampled clock, oven 16x, that is used in the receiver to check in the middle of each received by instead.
Then when a start bit is received the receiver will wait for roughly half a baud

\subsection{UART Component}



\subsection{Test Bench}
In order to test the UART module a simple test bench was created

\section{Original Design}
The original design, after correcting a few errors, did the following:

\section{UART Calculator}
This original design was heavily modified, with the original encoder being fully removed and the decoder being drastically changed, in order to meet the requirements for this assignment.
Figure \ref{fig:toplevel} shows this modified design.
As the UART components remains unchanged it will not be discussed again.
The overall design has three inputs: the clock, a reset line and the uart in, and only one output: the uart out.
Additionally there are four key generics: the clock rate, the baud rate, the maximum integer value, and the minimum integer value.
The first two values are used only by the UART component so that the baud pulses can be generated properly,
the integer range on the other will somewhat dynamically create version of the calculator cable of handling larger integers.
In practice the limit is -32767 to 32767, larger values may work in simulation but timing constraints prevent them from being implemented on physical hardware.


\subsection{ASCII Decoder}

\subsection{Calculator}
\subsection{ASCII Encoder}
\subsubsection{Sender}

\subsection{Test Bench}

\section{Implementation}
\subsection{Timing}

\pagebreak
\printbibliography

\end{document}